\documentclass{article}
\usepackage[T1]{fontenc}
\usepackage{lmodern}
\usepackage{graphicx}
\usepackage{amsmath,amssymb}
\usepackage[a4paper,margin=2cm]{geometry}
\usepackage[absolute,overlay]{textpos}
\setlength{\TPHorizModule}{1mm}
\setlength{\TPVertModule}{1mm}
\usepackage[indonesian]{babel}
\usepackage{hyperref}
\setlength{\emergencystretch}{2em}
% unit: Halaman 41

\begin{document}

% === Halaman 41 (python/native) ===

41

Super-enkripsi

• Menggabungkan cipher substitusi dengan cipher transposisi.

• Disebut juga product cipher

• Mula-mula pesan dienkripsi dengan cipher substitusi, selanjutnya hasilnya 

dienkripsi dengan cipher transposisi (atau sebaliknya).

   Contoh.  Plainteks hello world

✓Mula-mula dienkripsi dengan caesar cipher  (k = 3) menjadi KHOOR ZRUOG

✓kemudian hasil ini dienkripsi lagi dengan cipher transposisi  (k = 4):


KHOO


RZRU


OGZZ

• Cipher modern menggunakan konsep kombinasi cipher substitusi dan cipher 

transposisi, namun operasinya dibuat sekompleks mungkin

→ Cipherteks akhir adalah: KROHZGORZOUZ

\end{document}
